\chapter{Breast Abscess}
\index{Breast Abscess}

\section*{Definition and Overview}
A breast abscess is a localised collection of pus that forms within the breast tissue, almost always as a consequence of bacterial infection. It most commonly develops as a complication of \textbf{mastitis} -- inflammation and infection of the breast -- in breastfeeding women, in whom it typically arises when infection is not cleared quickly enough and pus accumulates in a walled-off pocket. It can also occur in women who are not breastfeeding, where it is usually related to inflammation of the ducts behind the nipple (periductal mastitis), and it occasionally affects men and infants. An abscess is more than an infection: because pus is enclosed in a cavity that antibiotics penetrate poorly, treatment almost always requires drainage of the pus in addition to antibiotics. Prompt recognition and management prevent tissue destruction, recurrence, and complications such as fistulae.

\section*{Epidemiology}
Breast abscesses are most common in women of childbearing age, particularly during the first few weeks to months after childbirth. Around 1 in 10 breastfeeding women develops mastitis, and approximately 5--11\% of those progress to an abscess. In non-breastfeeding women the condition is less common and is strongly associated with smoking, which is also the dominant risk factor for periductal mastitis; this form tends to affect slightly older women and frequently recurs, especially in those who continue to smoke. Risk is increased in women with diabetes, obesity, or a weakened immune system. The widespread use of early ultrasound-guided aspiration has reduced the need for open surgical drainage and improved cosmetic outcomes.

\section*{Aetiology and Pathophysiology}
The infection is almost always bacterial. In breastfeeding women the most common organism is \textit{Staphylococcus aureus}, including methicillin-resistant strains (MRSA); streptococci, enterococci, and anaerobes are less common causes. The bacteria usually enter through a cracked, fissured, or damaged nipple, and milk stasis -- breast milk that is not being removed effectively -- allows organisms to multiply in the stagnant milk. In non-breastfeeding women the infection is typically polymicrobial, with anaerobes flourishing in the inflamed ducts behind the nipple and producing the foul-smelling pus characteristic of periductal mastitis.

Risk factors include:

\begin{itemize}
    \item Cracked, sore, or damaged nipples
    \item Infrequent or ineffective feeding and consequent milk stasis
    \item A weakened immune system or chronic illness such as diabetes
    \item Smoking, especially in non-breastfeeding women
    \item A previous breast abscess or recurrent mastitis
    \item Skin conditions or piercings that breach the skin barrier
\end{itemize}

\section*{Clinical Features}
\begin{itemize}
    \item A tender, swollen, and often reddened area of the breast, frequently in the periphery of the gland
    \item A firm or fluctuant (soft and fluid-filled) lump that does not settle with standard mastitis treatment
    \item Breast pain, which may be severe and throbbing
    \item Fever, chills, and a general feeling of being unwell
    \item Pus discharging from the nipple or through a skin opening in some cases
    \item Symptoms that begin as mastitis but worsen or fail to improve within a day or two of treatment
\end{itemize}

\section*{Differential Diagnosis}
\begin{itemize}
    \item Mastitis without an abscess -- infection without a discrete collection of pus
    \item Inflammatory breast cancer, a rare and aggressive form of breast cancer that can mimic infection
    \item A breast cyst or galactocele (a milk-filled cyst)
    \item Fat necrosis, damage to fatty breast tissue often following injury
    \item Fibroadenoma or other benign breast lumps
    \item Duct ectasia with secondary infection
\end{itemize}

\section*{When to Seek Medical Care}
See a doctor promptly if you are breastfeeding and develop a painful, reddened breast that does not improve with feeding, massage, and rest within 24--48 hours, or if a firm lump persists despite treatment. A lump that fails to resolve like an infection should always be assessed, because breast cancer can occasionally masquerade as an abscess.

\textbf{Contact your local emergency services for:}
\begin{itemize}
    \item High fever, chills, or shaking
    \item Rapid breathing, a rapid heart rate, or feeling faint -- possible sepsis
    \item Confusion or severe drowsiness
    \item Rapidly spreading redness over the breast or chest wall
\end{itemize}

\section*{Diagnosis and Investigations}
\begin{itemize}
    \item \textbf{History and examination:} assessment of the feeding history, the onset and progression of symptoms, and breast examination to feel for a soft, fluctuant lump and to look for skin changes
    \item \textbf{Ultrasound:} the key investigation, which confirms whether an abscess is present, shows its size and depth, distinguishes it from mastitis, and guides drainage
    \item \textbf{Microbiology:} culture of pus or breast milk to identify the bacteria and guide antibiotic choice, particularly when infection is severe or recurrent
    \item \textbf{Blood tests:} full blood count and inflammatory markers where the person is systemically unwell
    \item \textbf{Biopsy:} if a lump does not behave like an infection or fails to resolve after treatment, tissue sampling to exclude breast cancer
\end{itemize}

\section*{Management}
\begin{itemize}
    \item \textbf{Antibiotics:} prescribed by a doctor and active against the likely bacteria, adjusted where culture results indicate; a full course should always be completed
    \item \textbf{Drainage:} most abscesses require removal of the pus, either by needle aspiration, usually repeated and performed under ultrasound guidance, or by a small surgical incision for larger or loculated collections
    \item \textbf{Continued feeding or expressing:} in breastfeeding women, keeping the breast drained is essential, and feeding from the affected breast is generally safe for the baby
    \item \textbf{Pain relief:} simple analgesics such as paracetamol and ibuprofen
    \item \textbf{Warm compresses and gentle massage} to encourage drainage and relieve discomfort
    \item \textbf{Follow-up:} review to confirm the infection has fully resolved, to repeat ultrasound if a collection persists, and to ensure no lump remains
    \item \textbf{Surgery for recurrent disease:} in non-breastfeeding women with periductal mastitis, excision of the affected duct (Hadfield's operation) may be offered
\end{itemize}

\section*{Complications}
\begin{itemize}
    \item Recurrence of the abscess, particularly in smokers and after inadequate drainage
    \item A mammary fistula -- an abnormal channel connecting a duct to the skin, which can discharge persistently
    \item Chronic or persistent infection
    \item Scarring, dimpling, and cosmetic change to the breast
    \item Difficulty continuing breastfeeding and premature weaning
    \item Rarely, systemic spread of infection causing sepsis
    \item Delay in diagnosing an underlying breast cancer that was mistaken for infection
\end{itemize}

\section*{Prognosis and Outcomes}
Most breast abscesses resolve completely with drainage and an appropriate course of antibiotics, and the breast usually heals well with minimal cosmetic impact, particularly when ultrasound-guided aspiration rather than open drainage is used. In breastfeeding women the outlook is very good, although the experience may make continued feeding more difficult and some women choose to wean. In non-breastfeeding women, especially smokers, recurrence is more common and the condition may require repeated aspiration, a prolonged antibiotic course, or definitive duct surgery.

\section*{Prevention and Self-Care}
\begin{itemize}
    \item Seek early treatment for cracked, sore, or damaged nipples
    \item Feed frequently and effectively to prevent milk stasis, and check the baby's latch and positioning
    \item Seek early medical review for mastitis rather than waiting for it to progress to an abscess
    \item Complete the full course of prescribed antibiotics and attend follow-up
    \item Stop smoking, which is the main modifiable risk factor in non-breastfeeding women
    \item Manage diabetes and other chronic conditions carefully
    \item See a doctor about any breast lump that does not disappear after the infection has settled
\end{itemize}

\section*{Sources and Further Reading}
The following organisations provide reliable, evidence-based information for patients, students, and clinicians.

\begin{itemize}
    \item NHS. \textit{Breast abscess and mastitis: symptoms, causes, and treatment.} https://www.nhs.uk/conditions/
    \item Mayo Clinic. \textit{Breast abscess and mastitis: patient education.} https://www.mayoclinic.org/diseases-conditions/
    \item MedlinePlus. \textit{Breast abscess: consumer health information.} https://medlineplus.gov/
    \item MSD Manual Professional Version. \textit{Mastitis and breast abscess: clinical guidance.} https://www.msdmanuals.com/
    \item National Institute for Health and Care Excellence. \textit{Clinical guidance on mastitis and breast infections.} https://www.nice.org.uk/
    \item World Health Organization. \textit{Breastfeeding guidance, including management of mastitis.} https://www.who.int/
\end{itemize}
